%%%%(c)
%%%%(c)  This file is a portion of the source for the textbook
%%%%(c)
%%%%(c)    Abstract Algebra: Theory and Applications
%%%%(c)    Copyright 1997 by Thomas W. Judson
%%%%(c)
%%%%(c)  See the file COPYING.txt for copying conditions
%%%%(c)
%%%%(c)
\documentclass[11pt]{book}

\usepackage{enumerate} %% CPT

\usepackage{graphicx}
\usepackage{amssymb}
\usepackage{epstopdf}
\usepackage{amsmath}
\usepackage{amscd}   %% redundant after tikz conversions?
\usepackage{multicol}
\usepackage{longtable}
\usepackage{imakeidx}
\usepackage{placeins}%  CPT
\usepackage{url}
\usepackage{dialogue}
\usepackage{hhline}

\usepackage{tikz}
\newcommand{\tikzpreface}[1]{\relax}

\usepackage{float}
\restylefloat{table}

\usepackage{hyperref}
\hypersetup{
  colorlinks   = true, %Colours links instead of ugly boxes
  urlcolor     = blue, %Colour for external hyperlinks
  linkcolor    = blue, %Colour of internal links
  citecolor   = red %Colour of citations
}

% Widows and orphans, from VCU edition
% R Beezer, 2010/06/14
\widowpenalty=10000
\clubpenalty=10000
\raggedbottom

\setlength{\fboxrule}{0.5pt}

\setlength{\parskip}{1ex plus 0.5ex minus 0.2ex} % Added to put spaces between paragraphs.

\newcommand{\sectionvideohref}[2][6mm]{\href{https://www.youtube.com/watch?v=#2}{\raisebox{-0.5mm}[0pt]{\includegraphics[height=#1]{images/video_camera}}}}

%%%%(c)
%%%%(c)  This file is a portion of the source for the textbook
%%%%(c)
%%%%(c)    Abstract Algebra: Theory and Applications
%%%%(c)    Copyright 1997 by Thomas W. Judson
%%%%(c)
%%%%(c)  See the file COPYING.txt for copying conditions
%%%%(c)
%%%%(c)
%%%%%%%%%%%%%%%%%%%%%%%%%%%%%%%%%%%%%%%%%%%%%%%%%%
%
%  Macros for algebra text
%
%%%%%%%%%%%%%%%%%%%%%%%%%%%%%%%%%%%%%%%%%%%%%%%%%%


%********theorems, propositions, lemmas, corollaries, proofs**

%% RAB, 2008/01/28
%%   Ditched [chapter] at end of three subsidiary environments
%%   The "[chapter]" was being printed

\newtheorem{theorem}{Theorem}[chapter]
  
\newtheorem{proposition}[theorem]{Proposition}
 
\newtheorem{lemma}[theorem]{Lemma}
 
\newtheorem{corollary}[theorem]{Corollary}
 
\newenvironment{proof}{\noindent {\sc
Proof.}}{\hspace{\fill} $\square$}


%********New commands******

\newcommand{\notdivide}{{\not{\mid}}} %Changed command from \notmid to avoid conflict with tikz - TWJ 5/6/2010

\newcommand{\notsubset}{\not\subset}

\newcommand{\lcm}{\operatorname{lcm}}

\newcommand{\Hom}{\operatorname{Hom}} %Added operator \Hom - TWJ 8/19/2010

\newcommand{\cis}{\operatorname{cis}}

\newcommand{\exrule}{\vspace{-0.3in} \noindent \rule{5.00in}{.5pt}}
		%********rule for exercise headings*******

\newcommand{\drawbox}{\raisebox{3pt}{\framebox[1.7in]{\hspace*{1in}}}}
		%********draws an empty box**********

\newcommand{\histbox}{\raisebox{3pt}{\hspace*{4pt}\framebox[.5in]{\hspace*{.25in}}}}
		%********draws an empty box**********


%********New fonts*************

\newfont{\bfi}{cmbxti10 scaled\magstephalf}  	%11pt bold italic

\newfont{\cmss}{cmss10}				%10pt sans serif

\newfont{\histf}{cmr10}				%10pt roman
						%Historical notes








%********Headings for historical notes*****


\newcommand{\histhead}{
\vspace{3ex} \noindent \drawbox \hfill \hspace*{4pt}
{\bfi Historical Note} \hfill \drawbox
\vspace{2ex}
}


%%%%%%%%%%%%%%%%%%%%%%%%%%%%%%%%%%%%%
%
%  ETP Modifications:
%	- Chapters start new left or right, not new right
%	- Float numbers are to be bold and punctuated with a period, not colon
%	- Allow larger floats on a text page
%	- Re-work the Chapter Opener
%
%%%%%%%%%%%%%%%%%%%%%%%%%%%%%%%%%%%%%


\makeatletter

\setcounter{topnumber}{2}	% Default parameter value from LaTeX book style
\def\topfraction{.7}		% Default parameter value from LaTeX book style
\setcounter{bottomnumber}{1}	% Default parameter value from LaTeX book style
\def\bottomfraction{.3}		% Default parameter value from LaTeX book style
\setcounter{totalnumber}{3}	% Default parameter value from LaTeX book style
\def\textfraction{.15}		% CHANGE: parameter was {.2}
\def\floatpagefraction{.7}	% CHANGE: parameter was {.5}




% Modified to remove the period after the chapter and section numbers
% in the running heads
\def\ps@headings{\let\@mkboth\markboth
\def\@oddfoot{}\def\@evenfoot{}\def\@evenhead{\rm \thepage\hfil \sl
\leftmark}\def\@oddhead{\hbox{}\sl \rightmark \hfil
\rm\thepage}\def\chaptermark##1{\markboth {\uppercase{\ifnum \c@secnumdepth
>\m@ne
 \@chapapp\ \thechapter \ \ \ \fi ##1}}{}}\def\sectionmark##1{\markright
{\uppercase{\ifnum \c@secnumdepth >\z@
 \thesection \ \ \ \fi ##1}}}}
\def\ps@myheadings{\let\@mkboth\@gobbletwo
\def\@oddhead{\hbox{}\sl\rightmark \hfil
\rm\thepage}\def\@oddfoot{}\def\@evenhead{\rm \thepage\hfil\sl\leftmark\hbox
{}}\def\@evenfoot{}\def\chaptermark##1{}\def\sectionmark##1{}%
\def\subsectionmark##1{}}



% Modified to embolden the figure (or table) number and punctuate with period
\long\def\@makecaption#1#2{
 \vskip 10pt
 \setbox\@tempboxa\hbox{{\bf #1.} #2}
 \ifdim \wd\@tempboxa >\hsize{\bf #1.} #2\par \else \hbox
to\hsize{\hfil\box\@tempboxa\hfil}
 \fi}


% RAB, 2010/06/17
% Fancy chapter headings for print versions
% Do nothing for tex4ht versions
\font\bigbolditalic=cmsl10 scaled\magstep5
\def\@makechapterhead#1{%\vspace*{50pt}
 { \parindent 0pt \centering% was\raggedright
 \ifnum \c@secnumdepth >\m@ne\rule{2in}{.5pt}\hfill\raisebox{-.1in}{\fbox{\fbox{\bigbolditalic\thechapter\/}}}\hfill\rule{2in}{.5pt}\par%

 \vskip 20pt \fi \Huge \bf #1\par%
 \nobreak \vskip 40pt \framebox[\hsize]{\hspace*{1in}}}%
 \vskip 36pt plus 12pt minus 6pt }%
%
\def\@makeschapterhead#1{%\vspace*{50pt}
 { \parindent 0pt \centering% was \raggedright
 \hrule height .5pt\vspace{40pt}%
 \huge \bf #1\par%
 \nobreak \vskip 40pt \framebox[\hsize]{\hspace*{1in}}}%
 \vskip 36pt plus 12pt minus 6pt }%
%
% \clearpage below was \cleardoublepage
\def\chapter{\clearpage \thispagestyle{plain} \global\@topnum\z@%
\@afterindentfalse \secdef\@chapter\@schapter}





\makeatother





% Per-chapter items
\newcommand{\chaplabel}{\relax}
\newcounter{example}[chapter]

% Usage: \chap{title}{label}
\newcommand{\chap}[2]{
     \setcounter{example}{0}%
     \renewcommand{\chaplabel}{#2}
     \chapter{#1}\label{#2}
}

% Examples as environment
\newenvironment{example}[1]
  {\refstepcounter{example}{\bigskip\noindent\bf Example~\arabic{example}.}\label{example:\chaplabel:#1}}
  {\hspace\fill$\blacklozenge$\par\medskip}

%%% Other environments (CPT, Dec 22 2010) %%%
\newenvironment{prop}[1]
  {\refstepcounter{example}{\bigskip\noindent\bf Proposition~\arabic{example}.}\label{proposition:\chaplabel:#1}}
  {\par\medskip}

\newenvironment{exercise}[1]
  {\refstepcounter{example}{\bigskip\noindent\bf Exercise~\arabic{example}.}\label{exercise:\chaplabel:#1}}
  {\hspace\fill$\lozenge$\par\medskip}

\newenvironment{defn}
  {\refstepcounter{example}{\bigskip\noindent\bf Definition~\arabic{example}.}}
  {\hspace\fill$\triangle$\par\medskip}

\newenvironment{rem}
  {\refstepcounter{example}{\bigskip\noindent\bf Remark~\arabic{example}.}}
  {\hspace\fill$\triangle$\par\medskip}

\newenvironment{notation}[1]
  {\refstepcounter{example}{\bigskip\noindent\bf Notation~\arabic{example}.}\label{notation:\chaplabel:#1}}
  {\hspace\fill$\triangle$\par\medskip}

%%% Legacy from P&C
\newenvironment{thm}
  {\refstepcounter{example}{\bigskip\noindent\bf Proposition~\arabic{example}.}}
  {\par\medskip}

\newenvironment{eg}
  {\refstepcounter{example}{\bigskip\noindent\bf Example~\arabic{example}.}}
  {\hspace\fill$\lozenge$\par\medskip}

\newenvironment{exer}
  {\refstepcounter{example}{\bigskip\noindent\bf Exercise~\arabic{example}.}}
  {\hspace\fill$\lozenge$\par\medskip}

\newenvironment{exers}
  {\bigskip {\bf Exercises} \medskip}
  {\hspace\fill$\lozenge$\par\medskip}

\newenvironment{summary}
  {\bigskip {\bf Summary} \medskip \begin{itemize}}
  {\end{itemize}\hspace\fill$\lozenge$\par\medskip}

\newenvironment{problist}
  {\bigskip {\bf Problem List} \medskip \begin{itemize}}
  {\end{itemize}\hspace\fill$\lozenge$\par\medskip}

\newenvironment{scratchwork}
  {\medskip \par}
  {\medskip \par}

\newenvironment{solution}
  {\bigskip {\bf Solution} \medskip}
  {\hspace\fill$\lozenge$\par\medskip}

\newenvironment{warn}
  {\bigskip \noindent{\bf Warning} \medskip}
  {\hspace\fill$\lozenge$\par\medskip}

{\newcommand{\chapterletter}[1]{#1}}

\renewenvironment{corollary}
  {\refstepcounter{example}{\bigskip\noindent\bf Corollary~\arabic{example}.}}
  {\par\medskip}


% New commands defined by CPT (legacy from P & C)
\newcommand{\var}[1]{\sf{#1}} 
\newcommand{\terminology}[1]{\textbf{\textit{{#1}}}}
\newcommand{\term}[1]{\terminology{#1}} 
\newcommand{\hint}[1]{\emph{#1}} 
\newcommand{\cref}[1]{\sf{\term{ #1}}} %%%% Cross-reference, legacy from P & C
\newcommand{\Cref}[1]{\sf{\term{ #1}}} %%%% Cross-reference, legacy from P & C
\newcommand{\real}{\mathbb{R}}
\newcommand{\rational}{\mathbb{Q}}
\newcommand{\integer}{\mathbb{Z}}
\newcommand{\univ}{\mathbb{U}}
\newcommand{\factoidbox}[1]{#1}
\newcommand{\eif}{\Rightarrow}
\newcommand{\eiff}{\Leftrightarrow}
\newcommand{\qedhere}{\Box}
\newcommand{\qed}{\Box}
\newcommand{\compose}{\circ}
\newcommand{\bigset}[1]{ \left\{ #1 \right\}}
\newcommand{\Id}{\sf{Id}}
\newcommand{\rel}{\sim}

\renewcommand{\ref}{\href}


\newcommand{\class}[1]{\overline{#1}}
\newcommand{\divides}{\mbox{|}}

% Commands from Hildenbrand, "Induction"

\newcommand{\RR}{\mathbb{R}}
\newcommand{\QQ}{\mathbb{Q}}
\newcommand{\NN}{\mathbb{N}}
\newcommand{\ZZ}{\mathbb{Z}}
\newcommand{\CC}{\mathbb{C}}

\newcommand{\Quote}[2]{\begin{quote}
\slshape #1 \hfill{\hbox{\slshape ---#2}}\end{quote}}

\newcommand{\seq}[1]{\langle #1\rangle} 
\newcommand{\aaa}{\langle a\rangle}
\newcommand{\bbb}{\langle b\rangle}
\newcommand{\cc}{\langle c\rangle}
\newcommand{\liminfty}{\lim\limits_{n\to\infty}}

\providecommand{\ul}[1]{\underline{#1}}

% define dummy command if \pts is not used
\providecommand{\pts}[1]{}

% example macro
%\newcommand{\example}[1]{\begin{quote}\small
%\textbf{Example: }#1\end{quote}} 

%%%%%%%%%%%%%%%%%%%%%%%%%%%%%%%%%%%%%%%%%%%%%%%%%%%%%%%%%%%%%%%%%%%%%%%%%
% redefine \implies

\renewcommand{\implies}{\Rightarrow}

\providecommand{\equivalent}{\Leftrightarrow}

\newcommand{\dsum}{\displaystyle\sum}
\newcommand{\dprod}{\displaystyle\prod}

%% End of Hildenbrand


%
\makeindex
%
\pagestyle{headings}

% Following is recommended from http://mintaka.sdsu.edu/GF/bibliog/latex/floats.html
% Alter some LaTeX defaults for better treatment of figures:
    % See p.105 of "TeX Unbound" for suggested values.
    % See pp. 199-200 of Lamport's "LaTeX" book for details.
    %   General parameters, for ALL pages:
    \renewcommand{\topfraction}{0.9}	% max fraction of floats at top
    \renewcommand{\bottomfraction}{0.8}	% max fraction of floats at bottom
    %   Parameters for TEXT pages (not float pages):
    \setcounter{topnumber}{2}
    \setcounter{bottomnumber}{2}
    \setcounter{totalnumber}{4}     % 2 may work better
    \setcounter{dbltopnumber}{2}    % for 2-column pages
    \renewcommand{\dbltopfraction}{0.9}	% fit big float above 2-col. text
    \renewcommand{\textfraction}{0.07}	% allow minimal text w. figs
    %   Parameters for FLOAT pages (not text pages):
    \renewcommand{\floatpagefraction}{0.7}	% require fuller float pages
	% N.B.: floatpagefraction MUST be less than topfraction !!
    \renewcommand{\dblfloatpagefraction}{0.7}	% require fuller float pages

	% remember to use [htp] or [htpb] for placement

\begin{document}
%
%
\chap{Practice Test Questions}{TestPractice2}


\begin{exercise}{}
Let $S$ be a set with three elements.
\begin{enumerate}[(a)]
\item
How many subsets does it have?
\item
How many proper subsets does it have?
\item
How many nonempty subsets does it have?
\item
How many nonempty proper subsets does it have?
\end{enumerate}
\end{exercise}

\begin{exercise}{7}
What is the smallest number of elements a set must have in order to have 20 proper subsets?
\end{exercise}

\begin{exercise}{12}
\begin{enumerate}[(a)]
\item
Find three sets $A_1, A_2, A_3$ which are \emph{not} disjoint such that $A_{1} \cup A_2 \cup A_3 = {\mathbb R}$ and $A_{1} \cap A_2 \cap A_3 = \emptyset$
\end{enumerate}
\end{exercise}

\begin{exercise}{20}
Let ${\mathbb N}$ be the universal set and suppose that
\begin{align*}
A = \{ x \in {\mathbb N} : x \text{ is divisible by 5}\} \\ 
B = \{ x \in {\mathbb N} : x \text{ is divisible by 7}\} \\ 
C = \{ x \in {\mathbb N} : x \text{ is greater than 50}\} \\
D = \{\text{the even natural numbers}\}
\end{align*} 

\noindent
Specify each of the following sets. You may specify a set either by describing a property, by enumerating the elements, or as one of the four sets $A, B, C, D$:
\begin{enumerate}[(a)]
\item
$(A \cap B) \setminus C$
\item
$A \cap B \cap C \cap D$
\item
$A \cup B \cup C \cup D$

\end{enumerate}
\end{exercise}

\begin{exercise}{cap_group}
Given a set $A$, let $G$ be the set of all subsets of $A$. 
\begin{enumerate}[(a)]
\item
Does the set $G$  with the operation $\cap$ have the closure property? \emph{Justify} your answer.
\item
Does the set $G$  with the operation $\cap$ have an identity? If so, what is it? Which part of  Proposition~\ref{proposition:Sets:sets_theorem_set_ops} enabled you to draw this conclusion?
\item
Is the operation $\cap$ defined on the set $G$ associative? Which part of  Proposition~\ref{proposition:Sets:sets_theorem_set_ops} enabled you to draw this conclusion?
\item
Is the operation $\cap$ defined on the set $G$ commutative? Which part of  Proposition~\ref{proposition:Sets:sets_theorem_set_ops} enabled you to draw this conclusion?
\item
Does each element of $G$ have a unique inverse under the operation $\cap$? If so, which part of  Proposition~\ref{proposition:Sets:sets_theorem_set_ops} enabled you to draw this conclusion? If not, provide a counterexample.
\item
Is the set $G$ a group under the $\cap$ operation?  \emph{Justify} your answer.
\end{enumerate}
\end{exercise} 

\begin{exercise}{}
List the elements of the set
$\{y,g,Y,G\} \times \{y,g,Y,G\}$
\end{exercise}

\begin{exercise}{funtable}
Here is a function~$f$ given by a table of values.

\begin{center}
\begin{tabular}{c|c}
$x$ & $f(x)$ \\ \hline

1 & 2 \\
2 & -2 \\
3 & -6 \\
4 & -10 \\
5 & -14 \\
\end{tabular}
\end{center}

\begin{enumerate}[(a)]
\item  \label{FunctionByTableEx-domain}
What is the domain of~$f$?
\item \label{FunctionByTableEx-range}
What is the range of~$f$? 
\item  \label{FunctionByTableEx-f(3)}
What is $f(3)$?
\item  \label{FunctionByTableEx-pairs}
Represent $f$ as a set of ordered pairs.
\item  \label{FunctionByTableEx-formula}
Find a formula to represent~$f$.
\hyperref[sec:Functions:Hints]{(*Hint*)}
\end{enumerate}
\end{exercise}

\begin{exercise}{ExOrderPairs}
Let  $A = \{ \var{a}, \var{b}, \var{c}, \var{d} \}$ and  $B = \{1,3,5,7,9\}$.
Which of the following sets of ordered pairs represent functions from~$A$ to~$B$?

\begin{multicols}{2}
\begin{enumerate}[a.]
\item \label{WhichAreFuncsEx-a1b3c5d7e9a11}
$\{ (\var{a}, 1), (\var{b}, 3), (\var{c}, 5), (\var{d}, 7), (\var{a},9) \}$
\item \label{WhichAreFuncsEx-a1b3c5e7}
$\{ (\var{a}, 1), (\var{b}, 3), (\var{c}, 5) \}$
\item \label{WhichAreFuncsEx-a1b1c1d1e1}
$\{ (\var{a}, 1), (\var{b}, 1), (\var{c}, 1), (\var{d}, 1) \}$
\end{enumerate}
\end{multicols}
\end{exercise}

\begin{exercise}{27}
 For the given sets $A$ and~$B$:
\begin{enumerate}[(a)] 
\item Write each function from~$A$ to~$B$ as a set of ordered pairs. (It turns out that if $|A| = m$
and $|B| = n$, then the number of functions from~$A$ to~$B$ is~$n^m$. Do you see why?)
\item Find the range of each function.
\end{enumerate}
\begin{enumerate}[i.]
\item \label{FunctionsChapExers-FindAll-ab,cde}
$A = \{\var{a},\var{b}\}$, $B = \{\var{c},\var{d},\var{e}\}$
%\item $A = \{\var{a},\var{b},\var{c}\}$, $B = \{\var{d},\var{e}\}$
\end{enumerate}
\end{exercise}

\begin{exercise}{11Exers}
Either prove the function is one-to-one, or prove that it is not.
\begin{multicols}{2}
 \begin{enumerate}[(a)]
\item \label{11Exers-formula-f}
 $f \colon {\mathbb R} \to {\mathbb R}$ defined by $f(x) = x^3 - x$.
\item \label{11Exers-formula-g}
 $g \colon {\mathbb R} \to {\mathbb R}$ defined by $g(x) = x^3$.
\item \label{modular_g}
 $g \colon {\mathbb Z}_6 \to {\mathbb Z}_6$ defined by $g(x)= (x \oplus 2) \odot 2$ .
\item \label{modular_g}
 $g \colon {\mathbb Z}_7 \to {\mathbb Z}_7$ defined by $g(x)= (x \oplus 2) \odot 2$ .
\end{enumerate}
\end{multicols}
\end{exercise}

\begin{exercise}{11Exers}
Either prove the function is onto, or prove that it is not.
\begin{multicols}{2}
 \begin{enumerate}[(a)]
\item \label{11Exers-formula-f}
 $f \colon {\mathbb R} \to {\mathbb R}$ defined by $f(x) = x^3 - x$.
\item \label{11Exers-formula-g}
 $g \colon {\mathbb R} \to {\mathbb R}$ defined by $g(x) = x^3$.
\item \label{modular_g}
 $g \colon {\mathbb Z}_6 \to {\mathbb Z}_6$ defined by $g(x)= (x \oplus 2) \odot 2$ .
\item \label{modular_g}
 $g \colon {\mathbb Z}_7 \to {\mathbb Z}_7$ defined by $g(x)= (x \oplus 2) \odot 2$ .
\end{enumerate}
\end{multicols}
\end{exercise}

\begin{exercise}{LinearWhenBijectionExer}
Let $a,b \in \mathbb{R}$, and define $f \colon \mathbb{R} \to \mathbb{R}$ by $f(x) = a x + b$. 
\begin{enumerate}[(a)]
\item \label{LinearWhenBijectionExer-not0}
Show that if $a \neq 0$, then $f$ is a bijection.
\item \label{LinearWhenBijectionExer-0}
Show that if $a = 0$, then $f$ is \emph{not} a bijection.
\end{enumerate}
\end{exercise}

\begin{exercise}{} 
For each function, either prove that it is a bijection, or prove that it is not.
\begin{enumerate}[(a)]
\item \label{modular9}
 $g \colon {\mathbb Z}_9 \to {\mathbb Z}_9$ defined by $g(x)= (x \odot 3) \oplus  (x \odot 4)$ .
\item \label{modular_m6}
 $g \colon {\mathbb Z}_7 \to {\mathbb Z}_7$ defined by $g(x) = (x \odot 4) \oplus (x \odot 4) $ .
 \end{enumerate}
\end{exercise}

\begin{exercise}{ComposeExers-form} 
 The formulas define functions $f$ and~$g$ from~$\mathbb{R}$ to~$\mathbb{R}$. Find formulas for $(f \compose g)(x)$ and $(g \compose f)(x)$.
\begin{enumerate}[(a)]
\item \label{ComposeExers-form-(ax+b)(cx+d)}
 $f(x) = ax + b$ and $g(x) = c x + d$ (where $a,b,c,d \in \real$)
\item \label{ComposeExers-form-(|x|)(-x)}
 $f(x) = |x|$ and $g(x) = -x$ 
\end{enumerate}
\end{exercise}

 \begin{exercise}{VerifyInverseExers}
 In each case, determine whether $g$ is an inverse of~$f$.
 \begin{enumerate}[(a)]
 \item \label{VerifyInverseExers-(x^2)}
$f \colon \real^+ \to \real^+$ is defined by $f(x) =2x^2$ and 
 \\ $g \colon \real^+ \to \real^+$ is defined by $g(y) = \sqrt{y}/2$.
 \item \label{VerifyInverseExers-(sqrt(x+1)-1)}
$f \colon \real^+ \to \real^+$ is defined by $f(x) = \sqrt{x+1} - 1$ and 
 \\ $g \colon \real^+ \to \real^+$ is defined by $g(y) = y^2 + 2y$.
 \end{enumerate}
 \end{exercise}

\begin{exer} \label{DrawBinRelExer}
Let $A =  \{-2,-1,0,1,2\}$ Draw a digraph for each of the following binary relations on~$A$: 
 \begin{enumerate}[(a)]
 \item \label{DrawBinRelExer-married}
 $ R_c = \{\, (x,y) \mid  (x-y)^2 < 2 \,\} .$
  \item \label{DrawBinRelExer-lived}
$ R_d = \{\, (x,y) \mid  x\equiv y \pmod{2} \,\} .$
 \end{enumerate}
 \end{exer}

\begin{exercise}{17}
For each of the following, explain your answers. 
\begin{enumerate}[(a)]
\item Define the relation $\rel$ on $\mathbb{Z}$ as follows: $ a \rel b$ iff $|a - b|< 4$. Is $\rel$ transitive? Is it reflexive? Is it symmetric?
\end{enumerate}
\end{exercise}

\begin{exer} \label{BinRelSomePropsEx}
Find binary relations on $\{1,2,3,4\}$ that meet each of the following conditions 
(Express each relation as a set of ordered pairs, and draw the corresponding digraph.)
\begin{enumerate}[(a)]
\item \label{BinRelSomePropsEx-transandsymm}
transitive and symmetric, but not reflexive.
\end{enumerate}
\end{exer}

\begin{exercise}{EquivClassEasyEx}
Let $B = \{1,2,3,4,5\}$ and 
	$$S = \left\{ (1,1),\, (1,4),\, (2,2),\, (2,3),\, (3.2),\, 
		(3,3),\, (4,1),\, (4,4),\, (5,5)
		 \right\} .$$
Assume (without proof) that $S$ is an equivalence relation on~$B$. Find the equivalence class of each element of~$B$.
\end{exercise}

\begin{exercise}{}
Let $C = \{1,2,3,4,5,6\}$ and define $\rel_C$ by 
\[ x \rel_C y \iff x + y \text{ is divisible by 3.} \]
Draw the arrow diagram for $\rel_C$, and prove that it is an equivalence relation.
\end{exercise}

\begin{exercise}{Mod3TablesEx}
Make tables that show the results of:
\begin{enumerate}[(a)]
\item \label{Mod3TablesEx-multiplication}
multiplication modulo~$5$.
\item \label{Mod3TablesEx-subtraction}
subtraction modulo~$5$ (For $\class{a} - \class{b}$,  put the result in row $\class{a}$ and column; $\class{b}$.)
\end{enumerate}
For both (a) and (b), all table entries should be  either $\class{0} \ldots \class{4}$.
\end{exercise}

\begin{exercise}{ModArithEx2}  
Find $x,y \in \integer_{51}$ such that $x \neq \class{0}$ and $y \neq \class{0}$, but $x \cdot y = \class{0}$.
\end{exercise}

\begin{exercise}{WellDefEx}
 Show that$f(x)=x^2$ provides a well-defined function from~$\integer_5$ to~$\integer_5$. That is, show that if $a,b \in \integer$, such that 
\[ [a]_5 = [b]_5, \text{ then } [ a^2]_7 = [ b^2]_7.\]
\end{exercise}

\begin{exercise}{57}
Show that there is a well-defined function 
$f \colon \integer_4 \to \integer_{12}$, given by \\
$ f \bigl( [a]_4 \bigr) = [a]_{12}$. 
That is, show that if $[a]_{12} = [b]_{12}$, then $f \bigl( [a]_4 \bigr) = f \bigl( [b]_4 \bigr)$.
\end{exercise}

\begin{exercise}{66}
Let $f\colon \ZZ_8 \to \ZZ_8$ be defined by $f(x) =  x^2 $. Define a relation $\sim$ on $\ZZ_8$ by: $n \sim m$ iff $f(n) = f(m)$.
\begin{enumerate}[(a)]
\item
Show that $\sim$ is an equivalence relation: that is, show that $\sim$ is reflexive, symmetric, and transitive.
\item
According to Proposition~\ref{EquivRel->Part}, this equivalence relation produces a partition on  $\ZZ_8$. List the sets in the partition.
\end{enumerate}
\end{exercise}

\begin{exercise}{16}
With reference to a hexagon with vertices labeled $A,B,C,D,E,F$ counterclockwise, for the symmetries $f$ and $g$:
\begin{enumerate}[(i)]
\item
Write the symmetries $f$ and $g$ in tableau form.
\item
Compute $f \compose g$ and $g \compose f$, expressing your answers in tableau form.
\item 
Describe the symmetries that correspond to $f \compose g$ and $g \compose f$, respectively.
\end{enumerate}
\medskip
$f=$rotation by $ 180^\circ, g=$reflection across the line $CF$
\end{exercise}

\bigskip
\begin{exercise}{}
\begin{enumerate}[(a)]
\item
Write the symmetries of a square (vertices labeled $A,B,C,D$) in tableau form.
\item
Write the Cayley table for the symmetries of a square. You may use letters to represent each symmetry. 
\item
List the inverses of each symmetry of the symmetries of a square.
\end{enumerate}
\end{exercise}


\end{document}
